\begin{claim}{D.2.38}{Proved}{The exact order of the finite fourth-stage source class}
Let $M=M(\mathbb Z/2,3)$, with bottom and top mod-two homology
classes $A_0,B_0$, and let
\[
 D=D_2^{(4)}M=S^3_+\wedge_{C_2}M^{\wedge2}.
 \tag{D.2.38a}
\]
Here the action on $S^3$ is antipodal and the action on the smash
square interchanges the factors.  The class
$Y_{10}=\lambda_3(A_0,B_0)\in H_{10}(D;\mathbb F_2)$
is spherical, and
\[
 \pi_{10}D\cong(\mathbb Z/4)^2\oplus\mathbb Z/2.
 \tag{D.2.38b}
\]
Every element with mod-two Hurewicz image $Y_{10}$ has exact
order four.  In a stable cell model described below, the last
attaching map has the form
\[
 2y+\zeta+2d n+e t,\qquad d,e\in\mathbb F_2,
 \tag{D.2.38c}
\]
where $y$ is a free degree-ten sphere and
$\pi_{10}R=\mathbb Z/2\{\zeta\}\oplus
\mathbb Z/4\{n\}\oplus\mathbb Z/2\{t\}$ is the group of
the other ten-dimensional summand.  The conclusions do not
require the values of $d,e$.  In particular the last cell cannot
be treated as a pure degree-two Moore attachment.
The sphere $S^{10}\to C_4M$ used below to realize $Y_{10}$
has zero image under $C_4M\to QM$.
\end{claim}
\begin{proof}
We work at the prime two.  The complex $D$ is five-connected and
eleven-dimensional.  Freudenthal therefore identifies its ordinary
$\pi_{10}$ with stable homotopy.  Its suspension spectrum is
two-primary: its integral homology, computed below, is two-primary,
and the finite simply connected complex becomes contractible
after inverting two.  Thus the calculation also determines the
unlocalized group in \textup{(D.2.38b)}.

\medskip\noindent
\textbf{The configuration cells and their operations.}
The filtered stable splitting identifies \textup{(D.2.38a)}
with the weight-two summand of $C_4M$; see
\cite[Theorem~1.1, p.~577, and Theorem~3.2, p.~580]{Snaith1974}.
The homology operations and their top-operation corrections used
here are \cite[Chapter~III, Theorems~1.2--1.4,
pp.~216--219]{I_CohenLadaMay1976}.  All unlabeled homology in
the proof has coefficients $\mathbb F_2$.
The complete table needed in degrees six through eleven is
\[
\begin{array}{c|c|c|c|c}
\text{degree}&\text{class}&\beta&\operatorname{Sq}_*^2&
 \operatorname{Sq}_*^4\\ \hline
6&P=A_0^2&0&0&0\\
7&R=A_0B_0&P&0&0\\
7&S=Q^4A_0&P&0&0\\
8&T=B_0^2&0&P&0\\
8&U=Q^5A_0&0&P&0\\
9&V=Q^5B_0&0&S&0\\
9&W=Q^6A_0&U&0&0\\
10&X=Q^6B_0&V&U&P\\
10&Y=\lambda_3(A_0,B_0)&0&0&0\\
11&Z=Q^7B_0&Y&W&S
\end{array}
\tag{D.2.38d}
\]
For completeness, the integral chain calculation uses the Moore
chains $a_3,b_4$, $\partial b=2a$, and the antipodal resolution
cells $e_i$, $0\leq i\leq3$:
\[
 \partial(e_i xy)=e_{i-1}\bigl(xy+(-1)^i\tau(xy)\bigr)
                  +(-1)^i e_i\partial(xy),
 \tag{D.2.38e}
\]
with the first term omitted for $i=0$ and
$\tau(aa)=-aa$, $\tau(ab)=ba$, $\tau(ba)=ab$, $\tau(bb)=bb$.
Cancelling the three unit chain pairs gives elementary boundaries
$2,4,2,2,2$ in successive dimensions
$7\to6,8\to7,9\to8,10\to9,11\to10$.
Thus $H_i(D;\mathbb Z)$ is respectively
$\mathbb Z/2,\mathbb Z/4,\mathbb Z/2,\mathbb Z/2,
\mathbb Z/2$ for $6\leq i\leq10$, and zero in degree eleven.
The unit cancellations can be made in a CW-spectrum model.
We use that stable model below; chain diagonalization alone is
not being used to assert an unstable wedge decomposition.

\medskip\noindent
\textbf{An actual spherical representative of $Y$.}
The generalized Browder map on the bottom sphere and the unit is
an actual map
\[
 \Sigma^3(S^3\wedge M)=M(\mathbb Z/2,9)\longrightarrow C_4M.
 \tag{D.2.38f}
\]
Its restriction to the bottom $S^9$ has fourfold adjunct the
Whitehead square of the bottom $S^7\to\Sigma^4M$.
The octonionic multiplication makes $S^7$ an $H$-space, so
$[\iota_7,\iota_7]=0$; naturality makes that restriction null.
Collapsing the null bottom gives an actual sphere $S^{10}\to C_4M$
whose Hurewicz image is $Y$.  The generalized product/adjunction
interface is \cite[Definitions~2.2--2.3 and
Theorem~2.4]{EH5Arkowitz1962}, and its homology identification is
the Browder operation cited above.  Projecting stably to $D$
retains this weight-two class.  The stable range for $D$ then
gives a sphere in its ordinary $\pi_{10}$ with the same image.

\medskip\noindent
\textbf{The stable eight- and nine-dimensional models.}
Choose integral lifts of the mod-two coordinates
$z_7=S$, $c_7=R+S$, and $T'_8=T+U$, with orientations
chosen in the chain reduction \textup{(D.2.38e)}.
Their cellular boundaries are
$\partial z=2P$, $\partial T'=4c$, $\partial U=0$.
Table \textup{(D.2.38d)} gives
$\operatorname{Sq}_*^2 U=P$ and
$\operatorname{Sq}_*^2 T'=0$.
The only possible degree-seven attaching torsion is the bottom
$\eta$, detected by this second square.  Therefore
\[
 D_8=C\vee M(\mathbb Z/4,7),\qquad
 C=\operatorname{Cof}\bigl((2,\eta_6):S^6\vee S^7\to S^6\bigr).
 \tag{D.2.38g}
\]
Let $i:S^6\to C$ and $q:C\to S^7\vee S^8$ be its structure maps.
The direct Chang-cofiber calculations of D.2.18--D.2.20,
shifted by two, give
\[
 \pi_8 C=\mathbb Z_{(2)}\{\gamma\}\oplus\mathbb Z/2\{\delta\},
 \qquad q\gamma=(0,2),\quad q\delta=(\eta_7,0),
 \qquad \gamma\eta_8=0.
 \tag{D.2.38h}
\]
These calculations use only the cofiber exact sequence and
the low sphere stems and Moore coextension identities in
\cite[Section~2, equations~(2.1)--(2.2), p.~149]{I_Oka1986}.
The degree-nine cells $W,V$ have boundaries $2U,0$ and
second squares $0,z$, respectively.  Their possible Moore-four
components are detected by the absent $c$ coefficient of the
second square.  Their attachments are consequently exactly
$\gamma,\delta$.  Put
\[
 C_b=C\cup_\gamma e^9_W\cup_\delta e^9_V,
 \qquad D_9=C_b\vee M(\mathbb Z/4,7).
 \tag{D.2.38i}
\]
This notation includes the stable unit cancellations of
\textup{(D.2.38e)}; no cancelled nine-cell is retained.

\medskip\noindent
\textbf{The component groups and the decisive composition.}
The same cofiber sequences give
\[
 \pi_9 C_b=\mathbb Z_{(2)}\{\kappa\}\oplus\mathbb Z/2\{v\},
 \quad h_{\mathbb Z}(\kappa)=2V,
 \quad v=i\nu_6,
 \tag{D.2.38j}
\]
where $\kappa$ is normalized to have fourth-square detector zero.
Its second-square detector is $U$.  This is the degree-two
pinch/Bockstein and Adem calculation used for $\kappa'$ in
D.2.20; adding $v$ changes exactly the fourth-square detector.
One degree higher the exact sequence for $C$ gives
\[
 \pi_{10}C=\mathbb Z/4\{\omega\},\qquad
 q\omega=(2\nu_7,\eta_8^2),\qquad
 \delta\eta_8^2=2\omega.
 \tag{D.2.38k}
\]
Since $\gamma\eta=0$ and $\delta\eta\ne0$, attachment of
the two nine-cells gives
\[
 0\longrightarrow\mathbb Z/2\{\bar\omega\}
 \longrightarrow\pi_{10}C_b
 \longrightarrow\mathbb Z/2\{\bar\zeta\}
 \longrightarrow0.
 \tag{D.2.38l}
\]
The right quotient is $\eta$ on the $W$-cell.  Its extension
need not be split at this point.

The specific composition needed below is
\[
 \kappa\eta_9=\bar\omega.
 \tag{D.2.38m}
\]
To check it, first attach only the $\delta$-cell.
In $C'=C\cup_\delta e^9_V$ the group $\pi_{10}C'$ is
$\mathbb Z/2\{\bar\omega\}$.  A lift $\kappa'$ of degree
two on $V$ has second-square detector $U$.
After collapsing the bottom $S^6$, its $U$-component is $\eta_8$;
its further $\eta$ therefore has the nonzero $\eta_8^2$ component
and is $\bar\omega$.  Attaching $\gamma$ preserves this subgroup
because $\gamma\eta^2=0$.  Normalizing by $v$ has no effect,
since $\nu\eta=0$.  The right quotient in
\textup{(D.2.38l)} is detected by the second square to $W$.

For the other wedge summand the Moore coextension sequence gives
\[
\begin{split}
 \pi_9 M(\mathbb Z/4,7)
   &=\mathbb Z/2\{i_7\eta_7^2\}
      \oplus\mathbb Z/2\{\widetilde\eta\},\\
 \pi_{10} M(\mathbb Z/4,7)
   &=\mathbb Z/4\{n=i_7\nu_7\}
      \oplus\mathbb Z/2\{t=\widetilde\eta\eta_9\}.
\end{split}
\tag{D.2.38n}
\]
Here $\widetilde\eta$ is the order-two Moore-four coextension
already used in D.2.19--D.2.22.  In particular
$v\eta_9=0$ and $i_7\eta_7^3=4i_7\nu_7=0$.

\medskip\noindent
\textbf{The two ten-cells and the spherical safeguard.}
The boundary and square vectors in \textup{(D.2.38d)} imply
\[
 a_X=\kappa+v+\epsilon i_7\eta_7^2,
 \qquad \epsilon\in\mathbb F_2.
 \tag{D.2.38o}
\]
Its Moore-four coextension component would have a nonzero
$T'$ second-square coefficient, which is absent.  The displayed
$i\eta^2$ coefficient is invisible to these primary detectors
and is deliberately retained.  For $a_Y$, zero integral boundary
and zero second/fourth squares leave only a possible $i\eta^2$.

The spherical class constructed in \textup{(D.2.38f)} forces
$a_Y=0$, for the following integral reason.  Its stable sphere
lifts to the ten-skeleton.  If its degrees on the two ten-cells
are $(a,b)$, the cofiber sequence requires
$a a_X+b a_Y=0$ in $\pi_9D_9$.
The free $\kappa$ coefficient of $a_X$ forces $a=0$.
Its Hurewicz image forces $b$ odd, and multiplication by such
$b$ is invertible on the remaining two-primary torsion.
Consequently $a_Y=0$ and there is a stable splitting
\[
 D_{10}=R\vee S^{10}_y,\qquad
 R=D_9\cup_{a_X}e^{10}_X.
 \tag{D.2.38p}
\]
This uses the free $\kappa$ coefficient to rule out a correction
by $X$; mod-two sphericity alone would not suffice.

In the cofiber sequence for $a_X$, the map
$\pi_9S^9\to\pi_9D_9$ is injective because of that same free
coefficient.  There is no right-hand kernel term in the
calculation of $\pi_{10}R$.
Equations \textup{(D.2.38m)} and \textup{(D.2.38n)} give
$a_X\eta_9=\bar\omega$, independently of $\epsilon$.
Killing the left subgroup of \textup{(D.2.38l)} therefore gives
the exact direct sum
\[
 \pi_{10}R=\mathbb Z/2\{\zeta\}\oplus
           \mathbb Z/4\{n\}\oplus\mathbb Z/2\{t\}.
 \tag{D.2.38q}
\]
The class $\zeta$ is the image of a lift of $\bar\zeta$.

\medskip\noindent
\textbf{The final attachment.}
Its integral boundary is $2Y$, so it is $2y+z$ with
$z\in\pi_{10}R$.
The second-square vector is $W$.  Among the torsion summands
in \textup{(D.2.38q)}, precisely $\zeta$ has that detector,
so its coefficient is one.
The fourth-square vector is $z_7=S$, with zero $c_7$ coefficient.
The class $\zeta$ comes from $C_b$, so has zero $c_7$ detector.
The class $t$ factors through $\eta$ on $S^9$; naturality from
its cone $C\eta_9$, which has no degree-seven cohomology,
also gives zero $c_7$ detector.
The $c_7$ fourth-square detector of $n=i_7\nu_7$ is one,
by the Hopf invariant of $\nu$.  Its coefficient must thus
be even.  This proves \textup{(D.2.38c)}.

Eliminating $\zeta$ from the resulting presentation gives
$(\mathbb Z/4)^2\oplus\mathbb Z/2$:
the relation $2\zeta=0$ becomes $4y=0$, independently of
$d,e$.  Moreover $2y$ is nonzero, since modulo
$2\pi_{10}R$ its value has the nonzero $\zeta$ coordinate.
Every representative of $Y$ is an odd multiple of $y$ plus
torsion from \textup{(D.2.38q)}.  Its double retains this
nonzero $\zeta$ coordinate, while its fourth multiple vanishes.
It therefore has exact order four, as asserted.

\medskip\noindent
\textbf{The image of the actual sphere factor in $QM$.}
The composite of \textup{(D.2.38f)} with $C_4M\to QM$ is
null: the generalized Browder product becomes null after
stabilization, equivalently its generalized Whitehead adjunct
has zero stable suspension.  The Moore Puppe sequence therefore
places the image of its sphere factor in $2\pi_{10}QM$.
This subgroup is zero in the present degree, as can be checked
without making a general inference about factors of null maps.
Write $M_2=\operatorname{Cof}(2:S\to S)$, so
$\pi_{10}QM=\pi_7M_2$.
The map $2\id_{M_2}$ restricts to zero on the bottom sphere,
so it factors through $q:M_2\to\Sigma S$.
The Moore exact sequence gives
$\pi_1M_2=\mathbb Z/2\{i\eta\}$; hence this factorization
is either zero or $i\eta q$.
For $x\in\pi_7M_2$, its pinch $qx$ lies in
$\pi_6S=\mathbb Z/2\{\nu^2\}$, and
$\eta\nu^2=(\eta\nu)\nu=0$.
The two sphere facts are
\cite[equation~(2.1)(e'),(g), p.~149]{I_Oka1986}.
Consequently $2x=0$ for every $x$, proving the asserted zero
image of this actual sphere $S^{10}\to C_4M\to QM$.
No ordinary map $D\to QM$ is inferred from the stable weight
projection in this argument.

This source calculation leaves the lower attaching coefficients
in \textup{(D.2.38o)} and \textup{(D.2.38c)} uncomputed.
It does not assert that the paired target functional vanishes
under every map $D\to V$, nor that a full fourth-suspension
extension exists or is impossible.
\end{proof}
